\section{Movimiento oscilatorio}
Fabriqué las figuras de Lissajous en tres dimensiones extendiendo el modelo de osciladores armónicos independientes a tres coordenadas. Cada coordenada $(x, y, z)$ obedece una ecuación de movimiento con su propia frecuencia, lo que permite generar curvas 3D al combinar razones enteras entre frecuencias.

\subsection{Modelo}
El sistema se representa como tres osciladores armónicos desacoplados:

\begin{equation}
\ddot{x} = -\frac{k_1}{m_1} x, \quad
\ddot{y} = -\frac{k_2}{m_2} y, \quad
\ddot{z} = -\frac{k_3}{m_3} z.
\end{equation}

Con $m_1 = m_2 = m_3 = 1$ y $k_1 = l^2$, $k_2 = n^2$, $k_3 = p^2$, las frecuencias quedan $\omega_x = l$, $\omega_y = n$, $\omega_z = p$. Las condiciones iniciales se fijan en $x_0 = y_0 = z_0 = 1$, $v_{x0} = 2.36$, $v_{y0} = v_{z0} = 0$.

\subsection{Casos y configuración}
Se evalúan siete ternas $(l, n, p)$ para generar familias distintas de curvas 3D:

\begin{equation}
\{(1,2,1), (1,3,1), (2,3,1), (3,4,1), (3,5,2), (4,5,3), (5,6,4)\}.
\end{equation}

El tiempo de integración es $t_\text{fin} = 50$ con paso $h = 0.01$. Para cada terna se almacena la trayectoria completa y se grafica en un mosaico $3\times3$ con \texttt{plot3}.

\subsection{Implementación numérica}
El avance temporal se realiza con el método de Euler-Cromer. Primero se actualizan las velocidades con las aceleraciones instantáneas y luego se actualizan las posiciones. Esta variante mejora la estabilidad frente al Euler simple cuando se trata de sistemas oscilatorios.

\inputminted[linenos, frame=lines, fontsize=\footnotesize,
    breaklines=true
]{matlab}{code/problem04.m}
\captionof{listing}{Simulación de Lissajous 3D usando Euler-Cromer y visualización en subplots.}
\label{lst:prob4_1}

\vspace{10pt}

En cada subplot se usa \texttt{axis equal} para mantener proporciones reales entre ejes y \texttt{view(30,30)} para una perspectiva uniforme. El resultado es un conjunto de curvas de Lissajous tridimensionales que muestran cómo las razones enteras entre frecuencias controlan la geometría del trazo.